\section{Directed error enclosures and conditional zero counts}
\label{sec:directed}

To evaluate \(G_N\), one needs \(\beta_1,\ldots,\beta_{N-1}\).
The amplitude in \(L_N\) also uses \(\beta_N\), which is the first
omitted link. Its sign is known even when its value has not been
determined. Define
\begin{equation}
 \begin{aligned}
 d_N(z)&=(-1)^Nz^{2N},\qquad
 \alpha_N=\frac{\prod_{j=1}^{N}\beta_j}{(2N)!},\\
 \overline\alpha_N&=
 \frac{\prod_{j=1}^{N-1}\beta_j}{(2N)!}
                      (B_{N-1}-\beta_{N-1}),
 \qquad 0<\alpha_N\le\overline\alpha_N.
 \end{aligned}
 \label{eq:amplitude-budget}
\end{equation}
At \(N=1\), use \(\beta_0=0\) and the empty product one.
For \(z\ne0\), dividing the exact error by its known direction gives
\begin{equation}
 \frac{F(z)}{d_N(z)}=w_N(z)+\alpha_NE_N(z),\qquad
 w_N(z)=\frac{G_N(z)}{d_N(z)}.
 \label{eq:normalized-error}
\end{equation}
A zero of \(F\) would require cancellation of the two terms on the
right. We now describe a region of possible values of the second
term, retaining its positive amplitude.

Write \(D(c,r)=\{v\in\mathbb C:|v-c|\le r\}\). At a point
in \(|\Im z|\le Y\le1\), put \(\varepsilon=\epsilon_N(|z|,Y)\).
The allowed normalized error lies in
\begin{equation}
 \begin{split}
 \mathcal U_N(z)&=\bigcup_{0\le a\le\overline\alpha_N}D(a,a\varepsilon)\\
 &=\operatorname{conv}\bigl(\{0\}\cup
 D(\overline\alpha_N,\overline\alpha_N\varepsilon)\bigr).
 \end{split}
 \label{eq:directed-set}
\end{equation}
The equality follows by writing each disk as a scaling, by
\(a/\overline\alpha_N\), of the last disk. The union is compact
and convex. Geometrically it is the hull of the origin and a disk
centered on the positive real axis; it is not a disk centered at zero.
Including amplitude zero will be essential for a contour homotopy,
although the actual amplitude is strictly positive.

\begin{proposition}[Pointwise consequences]
\label{prop:directed-exclusion}
At \(z\ne0\) in the stated strip, either of the following conditions
implies \(F(z)\ne0\):
\begin{align}
 \operatorname{dist}(-w_N(z),[0,\overline\alpha_N])
     &>\overline\alpha_N\epsilon_N(|z|,Y),
       \label{eq:capsule-test}\\
 \epsilon_N(|z|,Y)<1,\qquad \Re w_N(z)&\ge0.
       \label{eq:halfplane-test}
\end{align}
When \(\epsilon_N<1\), a zero necessarily satisfies
\begin{equation}
 0<|w_N(z)|\le\overline\alpha_N(1+\epsilon_N),\qquad
 |\arg(-w_N(z))|\le\arcsin\epsilon_N.
 \label{eq:necessary-sector}
\end{equation}
At \(z=0\), one uses \(F(0)=1\) directly.
\end{proposition}

\begin{proof}
The set \(\mathcal U_N\) is contained in the closed
\(\overline\alpha_N\varepsilon\)-neighborhood of the segment
\([0,\overline\alpha_N]\). Hence \eqref{eq:capsule-test}
prevents cancellation in \eqref{eq:normalized-error}.
If \(\varepsilon<1\), then
\(\Re(\alpha_NE_N)\ge\alpha_N(1-\varepsilon)>0\), proving
\eqref{eq:halfplane-test}. A point in a disk \(D(a,a\varepsilon)\)
with \(a>0\) has modulus at most \(a(1+\varepsilon)\), positive
real part, and argument at most \(\arcsin\varepsilon\) in
absolute value. The last assertion follows because the distance
from \(a\) to the line at that argument is \(a|\sin\theta|\).
Apply these facts to \(-w_N=\alpha_NE_N\).
\end{proof}

For an exact test of membership in the smaller set
\eqref{eq:directed-set}, a complex value \(v\) belongs to it if and
only if
\begin{equation}
 \min_{0\le a\le\overline\alpha_N}
 \bigl\{|v|^2-2a\Re v+(1-\varepsilon^2)a^2\bigr\}\le0.
 \label{eq:directed-membership}
\end{equation}
This is just \(|v-a|^2\le a^2\varepsilon^2\). When
\(\varepsilon<1\), its minimizer is
\(\Re v/(1-\varepsilon^2)\), restricted to the indicated segment.
When \(\varepsilon\ge1\), an endpoint minimizes the affine or
concave quadratic. Equality in this closed-set test gives no exclusion.

\subsection{From a boundary comparison to a zero count}

The argument principle counts zeros inside a positively oriented
simple closed contour \(\Gamma\) by
\((2\pi\ii)^{-1}\int_\Gamma f'/f\), provided \(f\) is holomorphic
inside and on the contour and nonzero on its boundary. It counts
multiplicities: a double zero contributes two. The count is unchanged
during a continuous holomorphic deformation that never vanishes on
the boundary \cite[\S9.8]{topping}.

\begin{theorem}[Conditional transfer of a model zero count]
\label{thm:contour}
Let \(\Gamma\) be a positively oriented piecewise \(C^1\) simple
closed contour in \(|\Im z|\le Y\le1\). Suppose
\begin{equation}
 \operatorname{dist}(-w_N(z),\mathcal U_N(z))>0
       \qquad(z\in\Gamma,\ z\ne0).
 \label{eq:contour-separation}
\end{equation}
Then \(F\) and \(G_N\) are nonzero on \(\Gamma\) and have the
same number of interior zeros, with multiplicities.
\end{theorem}

\begin{proof}
Set \(H_t=G_N+t(F-G_N)\), \(0\le t\le1\). For nonzero boundary
points its normalized correction is \(t\alpha_NE_N\), which
belongs to the closed set \(\mathcal U_N\). Condition
\eqref{eq:contour-separation} therefore keeps every \(H_t\)
nonzero there. If the contour passes through zero,
\(H_t(0)=1\) for all \(t\). Continuity on the compact set
\(\Gamma\times[0,1]\) gives a positive minimum modulus.
The argument-principle integral is consequently continuous in \(t\)
and integer-valued, hence constant. Its endpoints give the result.
\end{proof}

The separation concerns values in the complex value plane along a
contour in the \(z\)-plane. One must still produce the contour and
know the model count. For example, separation plus a zero model
count excludes zeros of \(F\) inside. Pointwise
\eqref{eq:halfplane-test} can hold even where \(G_N=0\); such a
point cannot satisfy \eqref{eq:contour-separation}, since
\(0\in\mathcal U_N\). Pointwise exclusion and count transfer
have different hypotheses.

\subsection{Using one more coefficient}

If \(\beta_N\) is known, so is the exact amplitude in \(L_N\).
The corrected model \(C_N=G_N+L_N\) then has the explicit bound
\begin{equation}
 |F(z)-C_N(z)|\le|L_N(z)|\epsilon_N(|z|,Y).
 \label{eq:corrected-disk}
\end{equation}
A strict boundary inequality
\(|C_N|>|L_N|\epsilon_N\) transfers its zero count to \(F\)
by Rouch\'e's theorem \cite[Theorem~9.21]{topping}. But those same
coefficients already specify \(G_{N+1}\), so a fair comparison is
with the next physical head. The leading errors obey
\begin{equation}
 \frac{L_{N+1}(z)}{L_N(z)}
   =-\frac{\beta_{N+1}z^2}{(2N+1)(2N+2)},
 \label{eq:leading-ratio}
\end{equation}
with the right side extending the ratio at zero. Since
\(\beta_{N+1}\le B_N-\beta_N\), for \(z\ne0\),
\begin{equation}
 \frac{|F-G_{N+1}|}{|L_N|}
 \le\frac{(B_N-\beta_N)|z|^2}{(2N+1)(2N+2)}
                 (1+\epsilon_{N+1}).
 \label{eq:next-head-budget}
\end{equation}
On ranges where the phase bounds stay bounded, this error bound is
\(O(|z|^2\log^2(N+1)/N^2)\); the displayed correction bound
in \eqref{eq:corrected-disk}, after division by \(|L_N|\), is
\(O(|z|^2\log^2(N+1)/N)\).

There is also a comparison of the actual errors at fixed degree:
\begin{equation}
 \lim_{z\to0,\ z\ne0}
 \frac{F(z)-G_N(z)-L_N(z)}{F(z)-G_{N+1}(z)}
 =\frac{\beta_N+\beta_{N+1}+2\sum_{j=1}^{N-1}\beta_j}
        {\beta_{N+1}}>1.
 \label{eq:local-cost-comparison}
\end{equation}
Indeed \(E_N(z)=1-\sigma_N^2z^2/2+O(z^4)\); substitute
\eqref{eq:error-moments} and \eqref{eq:leading-ratio}.
Both errors have order exactly \(2N+2\), and their quotient
extends holomorphically near zero. Thus the next head has smaller
actual error on a sufficiently small punctured disk for each fixed
\(N\). This proves no ordering of actual errors throughout the strip.
